\section{Related Work}
\label{sec:related_work}

\subsection{When Vision Meets Code}

A growing body of work explores programmatic representations as intermediate reasoning steps for visual understanding, moving beyond end-to-end neural prediction toward interpretable and executable scene hypotheses. The vision-as-inverse-graphics paradigm frames visual understanding as recovering the generative process behind an observed image or scene. VIGA~\cite{yin2026viga} instantiates this idea by iteratively writing, executing, rendering, comparing, and revising Blender Python code, while BlenderGym~\cite{gu2025blendergym} provides a systematic benchmark for graphics editing across placement, lighting, materials, blend shapes, and geometry. EZBlender~\cite{zhang2026ezblender} further decomposes complex editing instructions with a Plan-and-ReAct framework, and BlenderMCP~\cite{blendermcp2025} connects Blender to large language models through the Model Context Protocol as a practical agent interface.

Beyond 3D editing, code-based representations have also been used for spatial and physical reasoning. Thinking with Spatial Code~\cite{chen2026spatialcode} parses video into explicit 3D boxes and semantic variables for LLM reasoning, VisPhyWorld~\cite{visphyworld2026} asks models to generate executable simulator code from visual observations, and OmniLottie~\cite{yang2026omnilottie} shows that structured visual programs can represent temporally coherent animations. Closest in spirit to our representation, Code-as-Room~\cite{yang2026codeasroom} synthesizes executable Blender room scenes from top-down reference images through an agentic code-generation pipeline. BVB shares the view that code is a precise and verifiable medium for visual reasoning, but shifts the target from image-conditioned synthesis or static editing to reconstructing 3D scenes from video.

\subsection{Benchmarking Agents on Video}

Evaluating how well agents understand and act upon video content has become a central challenge, spanning passive video question answering, spatial reasoning, and active interaction. VSI-Bench~\cite{yang2024thinkinginspace} studies visual-spatial intelligence in real indoor videos and shows that spatial reasoning remains a major bottleneck for multimodal models, motivating structured intermediate representations for video understanding.
Other video-agent benchmarks focus on using video demonstrations or long-context video evidence to complete tasks in 2D digital environments. VideoGUI~\cite{lin2024videogui} evaluates GUI automation from instructional videos, VideoWebArena~\cite{jang2024videowebarena} evaluates web agents that require video understanding, and related GUI grounding work such as ScreenSpot-Pro~\cite{li2025screenspotpro} and GUI-Xplore~\cite{sun2025guixplore} targets high-resolution screen understanding and cross-application generalization. These benchmarks measure answer accuracy, action completion, or screen grounding. In contrast, BVB evaluates whether an agent can reconstruct the 3D content of a video inside a graphics engine, making preservation of layout, object relations, and QA-relevant spatial evidence testable through both rendering and code-level analysis.
